\newcommand{\F}{\mathbb{F}}
\newcommand{\C}{\mathbb{C}}
\newcommand{\an}{\mathscr{A}} 
\newcommand{\Q}{\mathbb{Q}}
\renewcommand{\hom}{\operatorname{Hom}}
\newcommand{\ehom}{\operatorname{End}}
\colorlet{gris}{black!60}

\-\vspace{-0.3cm}
  {\hrule \vspace{-0.15cm}
  \centering \textbf{Lista 2} \vspace{0.1cm}
  \hrule}

\vspace{0.2cm}
1. Sea \(G\) un grupo de orden \(68\). 

(a) Enuncie los teoremas de Sylow. 

\textcolor{gray}{\(\blacksquare\) \hyperlink{thm412}{4.12. (\emph{Sylow})} "(i) $p$-Sylow son conjugados e (ii) \(n_p \equiv 1 \ (\operatorname{mod} \ p)\) e \(n_p \mid |G|\)"; \(\blacksquare\) \hyperlink{thm414}{4.14.} "\(|G| = p^em : (p,m) = 1\) implica \(\forall P \leq G\) $p$-Sylow, \(|P| = p^e\)". } %e {\scriptsize \(\boxtimes\)} \hyperlink{coro415}{4.15.}. }

(b1) Muestre que \(\exists ! P \leq G\), \(17\)-Sylow. Describa \(G\) como \(P \rtimes Q\), definiendo \(Q\). A que grupos puede ser isomorfo \(Q\)?

\textcolor{gray}{
  \(|G|= 2^2 \cdot 17\), luego \(n_{17}\equiv 1 \ (\operatorname{mod} \ 17)\) e \(n_{17}\in\{1,\cancel{2},\cancel{4}\}\), por lo tanto \(n_{17} = 1\) e \(\exists ! P \lhd G : |P| = 17\). 
}

\textcolor{gray}{
  Sea \(Q\leq G\) subgrupo $2$-Sylow, como \((4,17)=1\) sigue que \(P \cap Q = \id\) e \(|PQ| = |P||Q| = 68 \ \leftrightharpoons \ PQ = G\). Luego, 
  \[G = P \rtimes Q : Q \cong \mathbb{Z}_4 \text{\ o \ } Q \cong \textbf{V}.\] 
}
\hspace{-0.1cm}(b2) Muestre que \(G\) es soluble.  

\newcommand{\Z}{\mathbb{Z}}
\textcolor{gray}{
  \(P \cong \Z_{17}\) e \(\nicefrac{G}{P} : |\nicefrac{G}{P}| = 2^2 \), por tanto ambos son abelianos (solubles), sigue que \(G\) también es soluble, \(\blacksquare\) \hyperlink{thm517}{5.17.}. 
}

(b3) Describa todos los \(\varphi \in \operatorname{Aut}(\Z_{17}) : |\varphi|= 2\) o \(|\varphi| = 4\).

\textcolor{gray}{
  \(\operatorname{Aut}(\Z_{17}) \cong \Z_{17}^\times \), vía \([\varphi_a : x \mapsto ax] \ \mapsto \ a\in \Z_{17}^\times\cong \Z_{16}\), como \((3,16) =1 \) para \(k= 0 ,\ldots, 15\) tenemos 
  \[\operatorname{Aut}(\Z_{17}) = \{\varphi_k: x \mapsto 3^k x  \ (\operatorname{mod} \ 17)\}.\]
  Ahí \(|\varphi_k| = |3^k| = \nicefrac{|3|}{(|3|,k)} = \nicefrac{16}{(16,k)}\). Finalmente, para \(|\varphi_k| = 2  \ \leftrightharpoons \ (16,k) = 8 \ \leftrightharpoons \ k=8\) e \(|\varphi_k| = 4 \ \leftrightharpoons \ (16, k) = 4 \ \leftrightharpoons \ k = 4\) o \( 12\).    \vspace{0.2cm}
  \hrule 
  \( \varphi_8 : x \mapsto 3^8 x = x\mapsto 16 x = x\mapsto -x \)  \hfill Orden $2$  \vspace{0.2cm}
  \hrule
  \( \varphi_4 : x \mapsto 3^4 x = x\mapsto 13 x  \) \hfill  Orden $4$  \vspace{0.2cm} \\ 
  \( \varphi_{12} : x \mapsto 3^{12} x = x\mapsto 4 x  \)  \vspace{0.2cm}
  \hrule
}
\vspace{0.2cm}
(b4) Muestre que si \(G\) no es abeliano, entonces no es único a menos de isomorfismo.   

\newcommand{\aut}{\operatorname{Aut}}
\textcolor{gray}{
  %Como vimos en (b2) podemos suponer \(P = \Z_{17}\), \vspace{0.2cm}\\  
  \(\star \) \ Sean \(Q = \Z_4 = \langle q\rangle \) e \(\theta_1 : Q \to \aut(P)\) tal que \(\theta_1: q \mapsto \varphi_{8}\), bien definido e no trivial, luego 
  \[G_1 := P \rtimes_{\theta_1} \Z_4 \text{ \ \ (no abeliano)}.\]
  \(\theta_1(q)(a) = q a q^{-1} \neq a \ \leftrightharpoons \ xa \neq ax, \ a \in P\). \vspace{0.2cm}\\ 
  \(\star \) \ Sean \(Q = \textbf{V} = \langle a,b\rangle \) e \(\theta_2 : Q \to \aut(P)\) tal que \(\theta_2 : a \mapsto \varphi_{8} \) e \(\theta_2 : b \mapsto \id_{\aut(P)}\), no trivial, luego   
  \[G_2 := P \rtimes_{\theta_2} \textbf{V} \text{ \ \ (no abeliano)}.\]
  Finalmente, \(G_1 \ncong G_2 \) pues \(\Z_4 \ncong \textbf{V}\), $2$-Sylow son conjugados (isomorfos) e isomorfismo lleva $2$-Sylow en $2$-Sylow. 
}
 
2. (a) Encuentre todos los $2$-Sylow de \(A_4\). 

\textcolor{gray}{
  \(|A_4| = 12 = 2^2 \cdot 3\), luego \(n_2 \equiv 1 \ (\operatorname{mod} \ 2)\) e \(n_2 \in \{1,3\}\)
  \[A_4 = \{\id, \underbrace{[(1 \ \ 2 \ \ 3)]}_{8 \text{ elem.}}, \underbrace{[(1 \ \ 2) (3 \ \ 4)]}_{3\text{ elem.}}\}\]
  Ninguno de estos elementos tiene orden $4$, por lo tanto los $2$-Sylow son isomorfos a \(\textbf{V}\) que consiste de identidad mas $3$ elementos de orden $2$, la única (\(n_2 = 1\)) posibilidad:  
  \[P = \{\id, [(1 \ \ 2)(3 \ \ 4)]\}\] 
}
(b) Encuentre todos los $2$-Sylow de \(S_4\). 

\textcolor{gray}{
  ssfd
}

(c) Demuestre todo \(G: |G| = p^2\), con \(p\) primo, es abeliano.